\subsection*{Chapter 32: Project: A Transformer Language Model from Scratch}

\paragraph{Exercise 32.1.} The pre-tokens are \verb*|x|, \verb*| |, \verb*| =|, \verb*| 10|,
\verb*|.|, \verb*|5|, \verb*|;|, \verb*| |, \verb*| y|. \verb*|x| matches the letters alternative.
Of the two spaces before \texttt{=}, the whitespace-with-lookahead alternative takes the first,
because the greedy \verb|\s+| must give back the last space so that the next character is not a
non-space. The second space then attaches to \texttt{=} through the ``optional space +
symbols'' alternative. \verb*| 10| is the digits alternative with its leading space. The decimal
point and \texttt{5} are separate pre-tokens: \texttt{.} is a symbol and \texttt{5} a digit run,
so numbers with decimals split into three pieces. \texttt{;} is a symbol; the next two spaces
behave as before, giving \verb*| | and \verb*| y|.

\paragraph{Exercise 32.2.} Snapping a guess forward to the next separator moves it by the
distance to the next document start. For exponential document lengths with mean 4\,KB, that
distance is itself exponential with mean 4\,KB (memorylessness), and the largest of $k$ such
shifts is about $4\,\text{KB}\cdot\ln k$: about 8\,KB for $k=8$ and 17\,KB for $k=64$. Next to chunk
sizes of 1.5\,GB or 190\,MB, the size imbalance is negligible ($<10^{-4}$ relative). What
imbalances the workers is \emph{processing time}: text varies in regex cost, processes are
scheduled unevenly, and there are stragglers. With $k=8$ chunks and 8 workers, the slowest chunk
sets the wall time. With $k=64$ chunks scheduled dynamically (\code{imap\_unordered} on a
\code{Pool}), fast workers take more chunks and the tail is at most about one small chunk; each
task also holds less memory.

\paragraph{Exercise 32.3.} Per layer: attention $4d^2=4.19$M, SwiGLU $3d\dff=8.65$M, so
12.85M; 16 layers give 205.5M. Untied embedding and head add $2Vd=65.5$M, for 271.1M in total
(plus $\sim$34K norm gains). Forward matmul FLOPs per token, excluding the embedding lookup:
$2\times(205.5\text{M}+32.8\text{M})\approx0.477$\,GFLOPs. The attention-score matmuls ($QK^\top$
and $PV$ over the full $T\times T$ square) add $4TdL$ per token: $0.134$\,GFLOPs at $T=2048$
(22\% of the total $0.611$) and $2.15$\,GFLOPs at $T=32{,}768$ (82\%). A causal kernel that
skips the masked half cuts the attention term roughly in half.

\paragraph{Exercise 32.4.} One feasible plan in 360 minutes: (1) a coarse sweep of 6 learning
rates on a $\sqrt{10}$ grid with 4-minute runs (24 min); (2) a fine sweep of 4 full-budget runs around
the best value, including one above the stability limit (80 min); (3) for each of 4 batch sizes,
3 learning rates with short runs (48 min), then one full-budget confirmation per batch size
(80 min). That is 232 minutes, leaving about two hours for a second seed of the most important
comparisons and for reruns. Every comparison uses the same token budget, so steps $=\lceil
\text{tokens}/(BT)\rceil$ and the schedule end change with $B$, and every run is logged with step
and wall-clock time.
